% FILE: artifacts/mixture-input-encoding_poster_ccn2026/parts/macros.tex
% ==================================================================================================
% Reusable toolkit: type sizes, math macros, layout grid, coordinate helpers,
% card / heading / icon / citation commands, and TikZ arrow / plane / shadow styles.
% Requires colors.tex and fonts.tex to be loaded first.

% ===== THE TYPE SYSTEM ============================================================================
% THE ONLY PLACE A POINT SIZE MAY BE CHOSEN. Never write \fontsize elsewhere -- not in a block, not in
% a figure, not inside a TikZ node. If a role needs a size that is not here, add a NAMED macro here.
%
% The poster is 147.32 x 172.72 cm at TRUE SCALE: 1 pt = 0.351 mm, so a point size IS a physical size
% (40 pt = a 14.1 mm em; 24 pt = 8.4 mm). Read every number below as a millimetre statement.
%
% TWO READABILITY FLOORS, read at different distances. Nothing may be set below its own.
%   TEXT   36 pt (12.7 mm em) -- anything read as language: prose, headings, captions, tags, notes,
%                               references, card content, INCLUDING text drawn inside a TikZ card.
%   FIGURE 24 pt ( 8.4 mm em) -- annotations inside a diagram: axis ticks, symbol labels, marks.
%
% INSIDE A FIGURE, TWO LEVELS AND ONLY TWO:
%   \figannotsize (30) -- the OBJECTS the figure is about: vectors, patterns, matrices.
%   \figticksize  (24) -- everything SUBORDINATE to them: axis names, coefficients, ticks.
% 30:24 is a step; 36:24 reads as an accident rather than a hierarchy. \figcapsize (36) is for WORDS
% about a figure -- its caption, a shared axis label, a colour key -- never for a symbol inside one.
%
% FONT SIZE IS ABSOLUTE INSIDE FIGURES, deliberately. A figure scales its COORDINATES (`scale=1.4')
% and never its fonts: no `transform shape', no \scalebox around a picture. Two consequences:
%   - Reusing a figure at another physical scale means RE-DECLARING these sizes (see supplement.tex).
%   - A label's CLEARANCE does not scale with it. Enlarging a picture moves a label further out in
%     coordinates while its glyph stays put, so it buys no room: raising a size forces a re-check of
%     every anchor in the figure, and collisions are fixed by MOVING THE ANCHOR, never by dropping
%     below the floor.
%
% THREE LEADINGS, NOT ONE:
%   PROSE   1.40x -- running text that wraps for several lines. Aerated, for reading at distance.
%   COMPACT 1.15x -- headings, captions, labels, notes: text that wraps to two or three lines and
%                    must stay a visual unit. A 1.4 leading pulls such a block apart.
%   FLAT    1.00x -- single-line nodes and icon glyphs, where leading is inert.
% When a gap is wanted that none of the three gives, ASK FOR THE GAP; do not invent a fourth ratio.
%
% Every step that can carry MATH needs a matching \DeclareMathSizes line in fonts.tex, or its math
% silently drifts out of the system.
%
% ----- The scale: 8 point sizes ------------------------------------------------------------------
% TEXT (floor 36)
\newcommand{\titlesize}{\fontsize{94}{108}\selectfont}      % poster title
\newcommand{\subtitlesize}{\fontsize{78}{90}\selectfont}    % title continuation
\newcommand{\headingsize}{\fontsize{52}{60}\selectfont}     % block headings, authors, display formulae
\newcommand{\cardtitlesize}{\fontsize{44}{50}\selectfont}   % card banner titles, heading kickers
\newcommand{\bodysize}{\fontsize{40}{56}\selectfont}        % prose that wraps            (PROSE 1.40)
\newcommand{\bodycompact}{\fontsize{40}{46}\selectfont}     % body size, tight            (COMPACT)
\newcommand{\notesize}{\fontsize{36}{42}\selectfont}        % citations, metadata, tags, notes -- FLOOR
% FIGURE (floor 24). \figcapsize shares \notesize's point size: a caption inside a figure and a note
% outside one are the same kind of small text, and one value keeps them from drifting apart.
\newcommand{\figcapsize}{\fontsize{36}{42}\selectfont}      % figure captions, shared axis labels, keys
\newcommand{\figannotsize}{\fontsize{30}{35}\selectfont}    % the figure's OBJECTS (vectors, matrices)
\newcommand{\figticksize}{\fontsize{24}{28}\selectfont}     % subordinate marks: ticks, axes -- FLOOR
% ICON GLYPHS -- FontAwesome in a disc. Not text: not read as language, exempt from both floors, and
% sized to the disc that holds them, never to the type scale. Named so they cannot drift either.
\newcommand{\iconsm}{\fontsize{20}{20}\selectfont}          % 11 mm disc  (\checkring)
\newcommand{\iconmd}{\fontsize{32}{32}\selectfont}          % 28 mm disc  (card banners, bullet marks)
\newcommand{\iconlg}{\fontsize{40}{40}\selectfont}          % 22-24 mm disc (step numbers, question mark)
\newcommand{\iconxl}{\fontsize{48}{48}\selectfont}          % QR legend glyph
\newcommand{\iconxxl}{\fontsize{57}{57}\selectfont}         % block-heading number disc

% ===== Math macros ================================================================================
% Vectors use \vec{} (arrow); bold is reserved for matrices (J, I, Sigma).
\newcommand{\vk}{\vec{\kappa}}
\newcommand{\vu}{\vec{u}}
\newcommand{\veta}{\vec{\eta}}
\newcommand{\vv}{\vec{v}}
\newcommand{\mI}{\mathbf{I}}
\newcommand{\mJ}{\mathbf{J}}
\newcommand{\vm}{\vec{m}}
\newcommand{\vn}{\vec{n}}
\newcommand{\Sig}{\boldsymbol{\Sigma}}
% In-math text label rendered in the math (serif) text font, so words match the symbols.
\newcommand{\mlab}[1]{\text{{\mathtextfont #1}}}
% An \underbracket's READING, pinned to \notesize INDEPENDENTLY of the formula's size: the gloss is not
% part of the statement, it is the words that READ the statement, and it must stay quiet. \mlab would
% follow the surrounding math size and grow with the formula. Use \brlab under any brace.
% NOTE: THE GLOSS, NOT THE ALGEBRA, SETS AN \underbracket's WIDTH -- it is the wider of the two.
\newcommand{\brlab}[1]{\text{{\mathtextfont\notesize #1}}}
% Inline math in a TikZ node that must read at prose size. A NAME ONLY: it must not transform anything.
% \DeclareMathSizes already reaches inside a TikZ node, so a \scalebox here would enlarge a second time.
\newcommand{\tmath}[1]{\ensuremath{#1}}
% multi-letter "mix" index set one math style smaller, to read like the 1/2 indices
\newcommand{\kmix}{\vk_{\scriptscriptstyle\mathrm{mix}}}
\newcommand{\umix}{\vu_{\scriptscriptstyle\mathrm{mix}}}

% ===== Layout grid (cm; origin at the page's south-west corner) ===================================
% Page (true CCN 2026 portrait size).
\def\PW{147.32}\def\PH{172.72}
% Columns (x). The two content columns stack their blocks INDEPENDENTLY (no rigid 2x2),
% which lets each side flow to fit its content.
% THE SIDE MARGIN IS THE ONE CHOSEN NUMBER HERE; the right edge and the full content width are
% ARITHMETIC ON IT AND \PW, not three numbers that have to be kept in step by hand.
\def\gMarg{5}                          % side margin, both sides
\edef\gXl{\gMarg}                      % left content edge
\edef\gXr{\fpeval{\PW-\gMarg}}         % right content edge            (142.32)
\def\gWl{65}           % the width \sbApprBody builds the GLOBAL move boxes at (block-approach.tex)
\edef\gWfull{\fpeval{\PW-2*\gMarg}}    % full content width (footer / supporting)   (137.32)

% ===== REGION MAP (cm; y = NORTH edge of a block, h = its height) =================================
% The two-question layout's own map (\yPhen / \hPhen, \yCol, \gXcolL / \gXcolR / \gGutC) lives in
% mixture-input-encoding-poster.tex, because it is specific to that layout. Only the two full-width anchors
% shared with the block files stay here:
\def\yHead{168.72}                           % header banner, full width. Its height is measured, not
                                             % set: read "[layout] header banner height" after any
                                             % header edit and re-derive the column tops.
\def\yFoot{12}                               % footer, full width

% Coordinate helper: \xy{x}{y} -> point at (x cm, y cm) from the page corner (O).
\newcommand{\xy}[2]{($(O)+(#1 cm,#2 cm)$)}

% ===== Card (panel) styles ========================================================================
% THE CARD GRAMMAR, and both halves of it are load-bearing:
%   NO BANNERS. A card's identity rides on a solid accent DISC + an accent-coloured TITLE on the plain
%   surface. A full-bleed accent band is the loudest device a card has, and it would be spent on the
%   card's NAME while the payoff below stays plain text. Do not reintroduce `colbacktitle'.
%   RADIUS IS A RATIO, NOT A NUMBER: ~4-5% of the card's own width (24 mm at 65 cm; 14 mm at 30 cm for
%   a nested card). Below ~3% a poster-scale arc reads as a chamfer. SCALE IT WITH THE CARD.
% CARD PADDING, HELD AS PLAIN cm NUMBERS. Two consumers need the same value in two forms: a tcolorbox
% key wants a dimension (`left=\gPadSideCm cm') and the region map wants a number it can do arithmetic
% on (the medallion's clearance against a question card's icon tile is derived from \gPadTopCm). One
% number serving both is what keeps the check measuring the point the layout actually places.
% TWO VERTICAL CONVENTIONS, deliberately distinct and each named: the plain `cardbase' pads 1.1 top and
% bottom, while a COLUMN card pads 1.3 / 1.4 because it carries a heading and a figure rather than a
% figure alone. Every measured box height in the poster is derived from one of these, so they are not
% interchangeable -- collapsing them means re-deriving \hBand and every row top.
\def\gPadSideCm{1.2}    % left / right, every card
\def\gPadBaseCm{1.1}    % top / bottom of a plain `cardbase' card
\def\gPadTopCm{1.3}     % top of a column card
\def\gPadBotCm{1.4}     % bottom of a column card
\tcbset{
  % shadow-free base: surface, rounded corners, padding.
  cardbase/.style={
    enhanced, arc=24mm, outer arc=24mm,
    colback=cardground, colframe=cardring, boxrule=0.8pt,
    boxsep=0pt, left=\gPadSideCm cm, right=\gPadSideCm cm,
    top=\gPadBaseCm cm, bottom=\gPadBaseCm cm,
    fontupper=\color{textprimary}\bodysize,
  },
  % ===== ONE SHADOW FOR EVERY CARD ON THE POSTER ==================================================
  % Every card glows in its identity colour through the SAME blur, cast and spread (`posterglow' below
  % -- the single copy of those numbers). Only the tint changes. TWO TRAPS, BOTH SILENT:
  %   (1) THE PRIMITIVE MATTERS AS MUCH AS THE NUMBERS. tcolorbox's `fuzzy shadow' is a stack of offset
  %       rounded rectangles; TikZ's `blur shadow' (shadows.blur) is a true gaussian blur, and giving a
  %       fuzzy shadow these numbers renders a stepped halo. Hence `blur shadow' on the FRAME PATH --
  %       the outermost path, drawn beneath the interior, so the shadow lands behind the whole card.
  %   (2) THE TINT MUST GO THROUGH `every shadow'. A `fill=' inside `blur shadow={...}' LEAKS onto the
  %       path it is attached to, and the card comes out flooded with the tint.
  %   #1 = ring colour (the frame path must keep filling itself with it)   #2 = tint of the glow
  cardglow/.style 2 args={
    frame style={fill=#1, posterglow, every shadow/.style={fill=#2, draw=none}}},
  % Section cards glow in their IDENTITY colour -- the heading, number disc and icon carry it -- and not
  % in their ring, which is a near-white hairline whose glow would be invisible. For a neutral cast:
  % {bordersubtle}{cardshadow}. The tint is `cardcast', so it follows the identity colour.
  card/.style={cardbase, cardglow={cardring}{cardcast}},
}

% A tcolorbox typeset directly inside a TikZ \node triggers an infinite NFSS font-selection loop under
% fontspec/LuaLaTeX. The fix is two-step: build the card into a savebox at document top level (OUTSIDE
% any tikzpicture), then place that box with \node{\usebox} inside the overlay picture.
% Build card CONTENT into <savebox> (call at document top level):
%   \buildcard{<\savebox>}{<width cm>}{<height cm>}{<tcb style>}{<content>}
\newcommand{\buildcard}[5]{%
  \sbox#1{\begin{tcolorbox}[#4, width=#2 cm, height=#3 cm]#5\end{tcolorbox}}%
}
% Build arbitrary (card-free) content into <savebox> at document top level:
%   \buildbox{<\savebox>}{<content>}
\newcommand{\buildbox}[2]{\sbox#1{#2}}

% Drop a pre-built box with its NORTH-WEST corner at (x,y) cm (call INSIDE the picture):
%   \dropbox{<\savebox>}{<x>}{<y>}
\newcommand{\dropbox}[3]{%
  \node[anchor=north west, inner sep=0pt] at \xy{#2}{#3} {\usebox#1};%
}
% As \dropbox but with an explicit anchor (default north west):
%   \dropboxa[<anchor>]{<\savebox>}{<x>}{<y>}
\newcommand{\dropboxa}[4][north west]{%
  \node[anchor=#1, inner sep=0pt] at \xy{#3}{#4} {\usebox#2};%
}

% Shared white space AFTER a body paragraph, before the figure row it introduces: a paragraph must not
% sit tight against the figure it announces.
% The space after a block HEADING is NOT set here -- it is \headbodygap (block-columns.tex), and no
% panel may add to it: heading, then content, nothing in between.
\newcommand{\paragap}{12mm}

% ===== Inline citation + supporting-detail reference tag ==========================================
% \litcite and \seebottom share the CHROME accent because they are the same kind of mark: a pointer
% AWAY from the poster (to the literature, to the supplement). Deliberately not two colours.
\newcommand{\litcite}[1]{{\color{pointermark}\notesize (#1)}}
% Small numbered badge matching the numbered points in supplement.tex. It is placed with \raisebox (a
% LaTeX box lift) around a picture whose baseline is pinned to the node CENTRE via a node anchor --
% neither depends on the surrounding picture's axes, so the badge aligns identically in ordinary text
% and inside a flipped-axis node (e.g. x=1cm y=-1cm).
\newcommand{\refnum}[1]{\raisebox{3.2mm}{\tikz[baseline=(rn.center)]{\node(rn)[fill=pointermark,
  text=white, rounded corners=3mm, minimum size=12mm, inner sep=0pt,
  font=\headingfont\notesize]{#1};}}}
% Pointer to numbered item(s) in the standalone supplement: the QR glyph followed by the item numbers,
% signalling "scan the header QR to reach supplemental item N". A BARE GLYPH SAYS ONLY `there is more
% somewhere'; the numbers are what make it a pointer, and they are the numbers the supplement's own
% \suphead badges carry -- keep the two in step when a supplement point is added or renumbered.
% Takes a COMMA LIST, so one macro serves one, two or three targets: \seebottom{4,5,6}.
% The glyph is set at \iconlg to match the heading it sits in; the badges bring their own size.
\newcommand{\seebottom}[1]{{\color{pointermark}%
  \mbox{{\iconlg\icQR}\;\foreach \sbn [count=\sbi] in {#1}{\ifnum\sbi>1\,\fi\refnum{\sbn}}}}}

% ===== THE SHARED 3-D CAMERA ======================================================================
% ONE elevation / azimuth for every tikz-3dplot figure that is NOT the mechanism tree or the supplement
% (those keep their own \mtcamtheta / \spcamtheta because they need a per-mode azimuth). \geometryplot,
% \inputplane / \responsedecomp and the two activity-space plots all read these, so the whole poster
% projects one scene from one viewpoint and a camera retune is a ONE-LINE edit here.
%   theta = ELEVATION (polar angle from the vertical z = n1). At 45 the view is OBLIQUE (the archive's
%   attitude, matching the mechanism tree): a tilted plane reads as tilted, not as a floor seen from above
%   (which is what the earlier low 29 gave). The span planes are kept OPEN not by lowering the camera but
%   by GENTLING the generators (z-share ~0.24): a plane whose normal sits only ~19 deg off vertical is
%   ~45% open even at this oblique view, while its residual still rises. phi = AZIMUTH, fixed so world-z
%   projects up, y right, x down-left (n1 up / n2 right / n3 receding), the archive's orientation.
\def\camtheta{57}\def\camphi{118}
% The ANSWER FACETS (\geometryplot) view the scene from FURTHER ABOVE (a smaller theta): this droops the
% floor axes more to the south, so the components plane reads as clearly distinct from the horizontal
% floor. Its plane is steepened to compensate (geometries.tex) so the plane's projected surface still
% matches the approach plots' look. phi is shared.
\def\facetcamtheta{38}
% ===== SHARED TikZ STYLES =========================================================================
% THE ONLY COPY OF EACH OF THESE NUMBERS. A figure that needs one of these devices takes the style and
% overrides only what is genuinely different (colour, and -- where noted -- weight); it never re-spells
% the arrow head or the shadow parameters. Retune here, or nowhere.
% ----- THE CARD GLOW'S THREE NUMBERS ---------------------------------------------------------------
% \providecommand, not \def, so that a document loading the colours first may pre-empt them.
% >>> WEIGHT MOVES WITH THE IDENTITY COLOUR <<< The glow is tinted with `cardcast', i.e. with the
% identity colour, and it is the accent's largest surface on the page. Under the deep petrol accent a
% strong cast reads as a bruise rather than a halo, so the glow is light and short: it reads as CONTACT
% rather than elevation. Retune these three with `chromecore'. Only `posterglow' reads them --
% `softshadow' is neutral depth INSIDE a figure.
\providecommand\gGlowDrop{-1.8mm}    % how far the glow falls below the card
\providecommand\gGlowBlur{4.5mm}     % its blur radius
\providecommand\gGlowOpacity{15}     % its strength, 0-100
\tikzset{
  % ---- shadows -------------------------------------------------------------------------------------
  % Deep glow: the poster's card / banner shadow, used by `cardglow' above and by the mechanism tree's
  % stage banners. The TINT is NOT here: it rides on `every shadow' at the call site (trap note above).
  posterglow/.style={blur shadow={shadow xshift=0pt, shadow yshift=\gGlowDrop,
    shadow blur radius=\gGlowBlur, shadow blur steps=12, shadow opacity=\gGlowOpacity}},
  % Shallow, neutral depth for objects INSIDE a figure (tree strips, option pills, taxonomy cells, leaf
  % plates). Lighter than `posterglow': they must not read as cards floating over the card holding them.
  softshadow/.style={blur shadow={shadow xshift=0pt, shadow yshift=-1.8mm, shadow blur radius=3.6mm,
    shadow blur steps=9, shadow opacity=34}},
  % ---- arrows / connectors ------------------------------------------------------------------------
  flowarrow/.style={-{Stealth[length=19pt,width=15pt]}, line width=2.2pt,
    rounded corners=8pt},
  % THE FRAME AXIS -- one weight and one head for every coordinate frame on the poster (the neuron
  % triples of \inputplane / \responsedecomp / \geometryplot, and the mechanism tree's intrinsic triple).
  % At 1.1 pt / 8 pt a frame prints as a hairline behind 3.4 pt vectors with invisible heads, so the
  % axes read as construction lines. The COLOUR is not here: each figure supplies its own `borderfunc' tint.
  frameaxis/.style={-{Stealth[length=15pt, width=11pt]}, line width=2.2pt},
  % A NAMED VECTOR -- the objects a figure is about. \geometryplot draws them at the declared 4 pt;
  % the Approach's two probe plots pass `line width=3.6pt' because they carry more arrows in less room.
  namedvec/.style={-{Stealth[length=16pt, width=12pt]}, line width=4pt},
  % THE ORTHOGONAL RESIDUAL (eta_perp): solid brown, one weight across the geometry figures.
  orthvec/.style={-{Stealth[length=13pt, width=10pt]}, line width=2.8pt, draw=outplane},
  % THE BIAS ARC -- the in-plane rotation off the linear sum, drawn concentric with it.
  biasarc/.style={-{Stealth[length=8pt]}, line width=1.4pt, draw=outplaneannot},
  % (The coordinate mark / dot lives further down as \dotonvec / \dotonaxis / \dotonplane, sized as a
  % fraction of each figure's axis length; see the block below the tikzset.)
  % THE LINEAR SUM (the proportional reference), one convention document-wide. Every geometry is
  % measured against it -- bias IS the deviation from it -- so it must be recognisable as the SAME
  % object in the taxonomy cells and in the tree. Absolute values, so `scale' cannot change its look.
  refvec/.style={-{Stealth[length=13pt, width=10pt]}, line width=2.8pt, draw=textsecondary,
    dash pattern=on 6pt off 3pt},
  % THE DROP LINE from a vector's tip (or its shadow) to the axis a coordinate is read on. ONE dash
  % geometry across both documents, so a reader crossing from poster to supplement meets the same mark.
  % Only the WEIGHT differs by figure family (`coorddrop' large, `coorddropsm' page-scale), because line
  % widths do not scale with a tikz `scale'. Override the colour with draw=.
  coorddrop/.style={line width=1.5pt, draw=texttertiary, opacity=0.9,
    dash pattern=on 3pt off 3pt, line cap=round},
  coorddropsm/.style={line width=0.9pt, draw=texttertiary, opacity=0.9,
    dash pattern=on 2.4pt off 2.4pt, line cap=round},
  % ---- the components plane -----------------------------------------------------------------------
  % THE HOUSE RULE FOR EVERY SPAN PLANE ON THE POSTER, stated once here: the plane IS the parallelogram
  % of its two generators -- 0 -> a -> a+b -> b -- so each arrow ENDS ON THE CORNER of the edge it
  % defines, and an inflated copy leaves every arrowhead floating short of its own corner. Anything
  % drawn inside the patch must keep BOTH coefficients < 1; the one licensed exception is
  % \geometryplot's unbiased cell, where the mixture IS the exact sum k1+k2 and lands on the far corner.
  % `compplane' IS "the plane of the components" here; two planes are told apart by ATTITUDE, never by
  % hue. Weight is overridden where a figure is drawn at a smaller scale.
  spanfill/.style={compplane, opacity=0.32},
  spanedge/.style={draw=compplane!75, line width=0.8pt},
}
% THE SHARED COORDINATE DOT -- SIZED FROM THE LINE WIDTH of the line it sits on, in THREE kinds, one per
% kind of line the project draws:
%   \dotonvec   -- on a RESPONSE / component vector (the thick rays)
%   \dotonaxis  -- on a FRAME AXIS or basis vector (incl. eigen-axes and eigenvectors)
%   \dotonplane -- a foot on a PLANE, marked by the dotted drop lines
% A dot reads as a bead threaded on its line, so its size follows THAT LINE'S WIDTH: a thick response
% vector carries a big bead, a hair-thin panel axis a small one. Because a figure drawn small also draws
% THINNER lines, this tracks figure size for free, with no per-figure scale knob. The law is affine --
% diameter = \dotbase + \dotgain x width -- not a bare multiple: the small \dotbase keeps a bead on a
% 0.7pt supplement axis legible instead of shrinking it to a speck, while \dotgain sets how fast the bead
% grows with the line. (Calibrated so the poster's 3.6pt response vectors carry the archive's ~8pt bead.)
% Each figure \def's only the widths it uses -- \lwvec / \lwaxis / \lwplane -- to its OWN line weights.
%   DRAWN AS NODES, not `circle' paths: inside `tdplot_main_coords' a `circle' is projected to an ELLIPSE,
%   whereas a circle NODE is placed at the projected point and drawn round in the page. No `transform
%   shape' (the size is absolute pt, like a line width). White halo under a coloured disc; halo/disc =
%   \dothaloratio, the rim being the gap between the two filled discs.
\def\dotbase{1.5}\def\dotgain{1.83}\def\dothaloratio{1.32}   % diameter (pt) = base + gain*linewidth
\def\lwvec{3.6}\def\lwaxis{2.0}\def\lwplane{1.5}             % defaults; each figure sets its actual weights
\newcommand{\dotcore}[3]{% #1 line width (pt number)  #2 fill colour  #3 (coordinate)
  \node[circle, fill=white, minimum size=\fpeval{\dothaloratio*(\dotbase+\dotgain*(#1))}pt, inner sep=0pt] at #3 {};%
  \node[circle, fill=#2,    minimum size=\fpeval{\dotbase+\dotgain*(#1)}pt,                 inner sep=0pt] at #3 {};}
\newcommand{\dotonvec}[2]{\dotcore{\lwvec}{#1}{#2}}
\newcommand{\dotonaxis}[2]{\dotcore{\lwaxis}{#1}{#2}}
\newcommand{\dotonplane}[2]{\dotcore{\lwplane}{#1}{#2}}

% ===== TAPERED / CURVED CONSEQUENCE ARROW =========================================================
% A FILLED, tapered, curved arrow (a narrow tail swelling to a wide head with a triangular point) --
% the "swoosh" connector between panels. No standard package draws a tapered CURVED body, so it is a
% filled bezier ribbon: the spine is a quadratic bend, the two edges are the spine offset by a width
% that grows from tail to neck, and a triangle caps the head. Widths are ABSOLUTE (pt), so the arrow
% keeps its weight regardless of the picture `scale' (like a line width or an arrowhead).
%   \taperarrow[<fill>][<extra path opts>]{<tail>}{<head>}{<bend>}{<tail hw>}{<neck hw>}{<ah len>}{<ah hw>}
%   bend = signed perpendicular offset of the spine's midpoint (0 = straight; +/- curves the swoosh).
%   The SECOND optional argument is a raw TikZ option list applied to the ribbon path together with the
%   fill -- pass `draw=<colour>, line width=<w>' for a border and a shadow key (e.g. `blur shadow={...}')
%   for a drop shadow. Empty by default, so bare `\taperarrow[<colour>]{...}' calls are unchanged.
\NewDocumentCommand{\taperarrow}{O{black} O{} m m m m m m m}{%
\begin{scope}
  \path
    let \p1 = ($#4-#3$), \n1 = {veclen(\x1,\y1)},
        \p2 = (\x1/\n1, \y1/\n1),            % unit direction (tail -> head)
        \p3 = (-\y1/\n1, \x1/\n1)            % unit perpendicular
    in
    coordinate (T@)  at #3
    coordinate (H@)  at #4
    coordinate (N@)  at ($#4-#8*(\p2)$)                      % neck (base of the arrowhead)
    coordinate (Bc@) at ($($#3!0.5!(N@)$)+#5*(\p3)$)         % bent spine mid
    coordinate (Tt@) at ($#3+#6*(\p3)$)  coordinate (Tb@) at ($#3-#6*(\p3)$)
    coordinate (Nt@) at ($(N@)+#7*(\p3)$) coordinate (Nb@) at ($(N@)-#7*(\p3)$)
    % control widths = the MEAN of the tail and neck half-widths, not the neck's: at the mean the four
    % control points of each edge are collinear when bend=0, so the taper is a straight ramp instead of
    % a curve that balloons out near the tail (visible only on a long arrow; short ones barely bent).
    coordinate (Ct@) at ($(Bc@)+0.5*#6*(\p3)+0.5*#7*(\p3)$)
    coordinate (Cb@) at ($(Bc@)-0.5*#6*(\p3)-0.5*#7*(\p3)$)
    coordinate (Aw1@) at ($(N@)+#9*(\p3)$) coordinate (Aw2@) at ($(N@)-#9*(\p3)$);
  \path[fill=#1, #2]
    (Tt@) .. controls (Ct@) and (Ct@) .. (Nt@) --
    (Aw1@) -- (H@) -- (Aw2@) --
    (Nb@) .. controls (Cb@) and (Cb@) .. (Tb@) --
    cycle;
\end{scope}}

% ===== THE POSTER'S CONNECTOR =====================================================================
% \connarrow{<tail>}{<tip>} -- ONE arrow shape for ONE relation ("leads to"), and the only \taperarrow
% call site the poster has. Every connector on the page runs through it: the two premise arrows into
% the medallion and the four question -> approach -> answer arrows down the columns. A second styling
% would assert a second kind of relation, which is exactly what the label system forbids.
% THE GEOMETRY IS TUNED FOR A SHORT ARROW. These runs are only ~2.6-3.7 cm, so the neck is thick and
% the head wide; a slender arrow at that length reads as a scratch. Straight (bend 0) because every
% run is axis-aligned.
% The border and shadow options go in \taperarrow's SECOND optional argument, and they are passed
% LITERALLY rather than through a macro -- pgfkeys must see the commas to split the list.
\newcommand{\connarrow}[2]{%
  \taperarrow[spinefill!45][draw=spineedge, line width=1.8pt, softshadow]%
    {#1}{#2}{0pt}{2pt}{18pt}{22pt}{44pt}}
% How far each end of a connector reaches INTO the two objects it joins, so it reads as one continuous
% link rather than as a floating glyph in a gap.
\def\connov{0.5}
